
\documentclass{article}
\usepackage{colortbl}
\usepackage{makecell}
\usepackage{multirow}
\usepackage{supertabular}

\begin{document}

\newcounter{utterance}

\centering \large Interaction Transcript for game `codenames', experiment `frequency\_low', episode 9 with reflexion{-}Qwen3.5{-}4B.
\vspace{24pt}

{ \footnotesize  \setcounter{utterance}{1}
\setlength{\tabcolsep}{0pt}
\begin{supertabular}{c@{$\;$}|p{.15\linewidth}@{}p{.15\linewidth}p{.15\linewidth}p{.15\linewidth}p{.15\linewidth}p{.15\linewidth}}
    \# & $\;$A & \multicolumn{4}{c}{Game Master} & $\;\:$B\\
    \hline

    \theutterance \stepcounter{utterance}  
    & & & \multicolumn{2}{p{0.3\linewidth}}{
        \cellcolor[rgb]{0.9,0.9,0.9}{
            \makecell[{{p{\linewidth}}}]{
                \texttt{\tiny{[GM$|$GM]}}
                {'hidden': {'team': ['backbreaker', 'nearness'], 'innocent': ['pleasantness', 'tonsillitis', 'defenseman', 'fortunetelling', 'secession'], 'opponent': [], 'assassin': ['shirtfront']}, 'revealed': {'team': {'team': [], 'innocent': [], 'opponent': [], 'assassin': []}, 'opponent': {'team': [], 'innocent': [], 'opponent': [], 'assassin': []}}}
            }
        }
    }
    & & \\ \\

    \theutterance \stepcounter{utterance}  
    & & \multicolumn{4}{p{0.6\linewidth}}{
        \cellcolor[rgb]{0.9,0.9,0.9}{
            \makecell[{{p{\linewidth}}}]{
                \texttt{\tiny{[P1$\langle$GM]}}
                \texttt{You have attempted this same task before and did not succeed. These are the plans you wrote after those attempts.} \\
\texttt{Trial \#1: New plan: The last attempt was aborted because the clue word "CLOSE" contained the substring "SE" which is part of the distractor word "secession", violating the rule that the clue cannot contain parts of the word lists. Next time, I will select a clue word that is semantically related to the target words but does not contain any substrings from the target, opponent, or distractor lists. For example, I could use "NEAR" for "nearness" or "SMALL" for "narrowness", ensuring the word is completely independent of the provided lists.} \\
\\ 
\texttt{Trial \#2: New plan: The last attempt was aborted because the clue word "NEAR" is a substring of the target word "nearness", violating the rule that the clue cannot contain parts of the target words. Next time, I will select a clue word that is semantically related to the target words but is completely independent, containing no letters or substrings from any of the provided lists (target, opponent, distractor, or kill). For example, I could use "PROXIMITY" for "nearness" or "HEIGHT" for "loftiness", ensuring the clue word has zero overlap with the provided vocabulary.} \\
\\ 
\texttt{Trial \#3: New plan: I failed because I chose "LOW" for "backbreaker," which is semantically related but likely too ambiguous given the opponent's remaining word "hotelkeeper" (which also relates to low rooms/levels) and the risk of guessing a distractor. I should have selected a more specific or unique association for "backbreaker" that distinguishes it from "nearness" and avoids overlap with opponent/distractor words. For the final turn, I will use a clue like "PAIN" or "EXHAUSTION" for "backbreaker" (focusing on the physical toll) and "DISTANCE" or "ADJACENCY" for "nearness," ensuring the clue words are completely independent of all provided lists and highly specific to the target meanings to force a correct guess without triggering opponent or distractor words.} \\
\\ 
\\ 
\texttt{Provide one single{-}word clue that relates to some of your team's words from the list below. You can choose to how many and to which words your clue relates to.} \\
\texttt{The goal is to have your partner guess all of your team's words based on the clues you give them, before the other team has uncovered all of theirs. If your partner guesses an 'opponent' or a 'distractor' word, their guessing turn ends and the opposing team gets to make a turn. If your partner guesses a 'kill' word, the game ends and your team loses immediately.} \\
\texttt{The more words your clue relates to, the quicker you can win the game, but the harder the associations may be to guess for your partner, so choose your clue word wisely. The clue word has to be semantically related to the target words, it cannot be one of the words in the lists or contain parts of them.} \\
\\ 
\texttt{Always give your single{-}word clue and your comma{-}separated list of related target words in the following format and make your answers as short as possible, never include any other text than is required in this form:} \\
\\ 
\texttt{CLUE: <WORD>} \\
\texttt{TARGETS: <TARGETS>} \\
\\ 
\texttt{Your team words are: backbreaker, nearness.} \\
\texttt{Your opponent's team words are: .} \\
\texttt{Distractor words are: pleasantness, tonsillitis, defenseman, fortunetelling, secession.} \\
\texttt{The kill words are: shirtfront.} \\
            }
        }
    }
    & \\ \\

    \theutterance \stepcounter{utterance}  
    & \multicolumn{4}{p{0.6\linewidth}}{
        \cellcolor[rgb]{0.9,0.9,0.9}{
            \makecell[{{p{\linewidth}}}]{
                \texttt{\tiny{[P1$\rangle$GM]}}
                \texttt{CLUE: PAIN} \\
\texttt{TARGETS: backbreaker} \\
            }
        }
    }
    & & \\ \\

    \theutterance \stepcounter{utterance}  
    & & & \multicolumn{2}{p{0.3\linewidth}}{
        \cellcolor[rgb]{0.9,0.9,0.9}{
            \makecell[{{p{\linewidth}}}]{
                \texttt{\tiny{[GM$|$GM]}}
                \texttt{pain} \\
            }
        }
    }
    & & \\ \\

    \theutterance \stepcounter{utterance}  
    & & & \multicolumn{2}{p{0.3\linewidth}}{
        \cellcolor[rgb]{0.9,0.9,0.9}{
            \makecell[{{p{\linewidth}}}]{
                \texttt{\tiny{[GM$|$GM]}}
                ['backbreaker']
            }
        }
    }
    & & \\ \\

    \theutterance \stepcounter{utterance}  
    & & & \multicolumn{2}{p{0.3\linewidth}}{
        \cellcolor[rgb]{0.9,0.9,0.9}{
            \makecell[{{p{\linewidth}}}]{
                \texttt{\tiny{[GM$|$GM]}}
                {'word': 'backbreaker', 'assignment': 'team'}
            }
        }
    }
    & & \\ \\

    \theutterance \stepcounter{utterance}  
    & & & \multicolumn{2}{p{0.3\linewidth}}{
        \cellcolor[rgb]{0.9,0.9,0.9}{
            \makecell[{{p{\linewidth}}}]{
                \texttt{\tiny{[GM$|$GM]}}
                \texttt{opponent has won} \\
            }
        }
    }
    & & \\ \\

    \theutterance \stepcounter{utterance}  
    & \multicolumn{4}{p{0.6\linewidth}}{
        \cellcolor[rgb]{0.95,0.95,0.95}{
            \makecell[{{p{\linewidth}}}]{
                \texttt{\tiny{[P2$\rightarrow$P1]}}
                \texttt{Post{-}episode Reflexion (not sent to GM):} \\
\texttt{New plan: I failed because "PAIN" is too generic and likely overlaps with distractor words like "tonsillitis" (which also relates to pain), causing the partner to guess a distractor instead of the target. Next time, I will select a clue word that is semantically related to "backbreaker" but has zero semantic overlap with any distractor words (especially "tonsillitis") and is distinct from "nearness." For "backbreaker," I will use a specific physical action or consequence like "STRETCH" or "INJURY" that is less likely to be confused with medical conditions, ensuring the partner guesses the exact target word without triggering a kill or distractor turn.} \\
            }
        }
    }
    & & \\ \\

\end{supertabular}
}

\end{document}
